\documentclass[12pt]{article}
\usepackage[utf8]{inputenc}
\usepackage[T1]{fontenc}
\usepackage{hyperref}
\title{West Albay Bikol Language Resources for K-12 Curriculum Development}
\author{Compiled by Claude Sonnet 5.5 (claude-sonnet-5-5) for Andr\'as Kornai}
\date{}
\begin{document}
\maketitle

\section{Introduction}
West Albay Bikol (ISO 639-3 \texttt{fbl}) is a Bikol variety of Albay province, Philippines, spoken in towns between the Rinconada area of Camarines Sur and Legazpi City, each with a noticeably different dialect. Searches found one grammatical sketch of the language, general Bikol reference works, and material on Bikol mother-tongue education in the region. Nothing specific to West Albay for schools was found.

\section{Primary Educational Resources}
\begin{itemize}
\item DepEd LRMDS: ``MTB-MLE Grade 1: Bikol Week 30 Teacher's Guide'' (Bikol, not specific to West Albay): \url{https://lrmds.deped.gov.ph/detail/3961}
\item DepEd Albay district memorandum on a training in Tiwi, Albay, to adapt Grade 1 Central Bikol materials into local mother tongues (Scribd copy): \url{https://www.scribd.com/doc/160685911/Training-Matrix}
\end{itemize}

\section{Linguistic and Reference Materials}
\begin{itemize}
\item Nhia Borja, 2024 CUNY thesis containing a grammatical sketch of West Albay Bikol, based on narratives from Balugo and Polangui speakers: \url{https://academicworks.cuny.edu/gc_etds/5781}
\item Bikol Dictionary (University of Hawai'i Press), based on the Naga standard rather than West Albay: \url{https://manifold.uhpress.hawaii.edu/projects/bikol-dictionary}
\item C. McFarland, \emph{The Dialects of the Bikol Area} (Yale dissertation, 1974), bibliographic record: \url{https://glottolog.org/resource/reference/id/61602}
\item Bikol phrasebook (Wikivoyage), covering Naga and Legazpi varieties: \url{https://en.wikivoyage.org/wiki/Bikol_phrasebook}
\end{itemize}

\section{Religious and Community Resources}
None found in these searches.

\section{Digital and Audio Resources}
None found for West Albay specifically.

\section{Curriculum Development Context}
RA 10533 requires K--3 instruction and materials in the learners' regional or native language. Bicol University studies report that Bikol Naga was assigned as the regional MTB-MLE language, that Central Bikol (Legazpi) materials are in use, and that teachers adapt them to local varieties:
\begin{itemize}
\item Bicol University R\&D Journal article (de la Torre): \url{https://systems.bicol-u.edu.ph/journal/assets/journal_pdf/1_de%20la%20Torre_1-12.pdf}
\item Bicol University R\&D Journal article (Lorenzana): \url{https://systems.bicol-u.edu.ph/journal/assets/journal_pdf/4%20Lorenzana_49-61.pdf}
\end{itemize}
A West Albay curriculum would have to build on the Borja grammatical sketch and adapt Central Bikol materials, as no West Albay-specific school materials were found.

\end{document}